\documentclass[12pt, a4paper]{article}

% -------- Packages --------
\usepackage[margin=2.5cm]{geometry}
\usepackage{cite}
\usepackage{amsmath,amssymb,amsfonts}
\usepackage{graphicx}
\usepackage{textcomp}
\usepackage{xcolor}
\usepackage{booktabs}
\usepackage{multirow}
\usepackage{array}
\usepackage{url}
\usepackage{hyperref}
\usepackage{setspace}
\usepackage{parskip}
\usepackage{titlesec}
\usepackage{caption}
\usepackage{subcaption}
\usepackage{enumitem}
\usepackage{fancyhdr}
\usepackage{abstract}
\usepackage{verbatim}

% -------- Figure path --------
\graphicspath{{./}}

\onehalfspacing

% -------- Header / Footer --------
\pagestyle{fancy}
\fancyhf{}
\rhead{\small Flight Fare Prediction}
\lhead{\small National Economics University}
\cfoot{\thepage}
\renewcommand{\headrulewidth}{0.4pt}

% -------- Section formatting --------
\titleformat{\section}{\large\bfseries}{\thesection.}{0.6em}{}
\titleformat{\subsection}{\normalsize\bfseries}{\thesubsection.}{0.5em}{}
\titleformat{\subsubsection}{\normalsize\itshape}{\thesubsubsection.}{0.5em}{}

% -------- Hyperref colours --------
\hypersetup{colorlinks=true, linkcolor=black, citecolor=black, urlcolor=blue}

% ===========================================================================
\begin{document}
% ===========================================================================

\begin{titlepage}
\centering
\vspace*{2cm}
{\Large\bfseries Flight Fare Prediction Using Machine Learning and\\[6pt]
Time Series Analysis: A Hybrid Approach for the\\[6pt]
Indian Domestic Aviation Market\par}
\vspace{1.5cm}
{\large [Author Name]\par}
{\normalsize Student ID: [Student ID]\par}
\vspace{0.5cm}
{\normalsize Faculty of Mathematical Economics\\
National Economics University\\
Hanoi, Vietnam\par}
\vspace{0.5cm}
{\normalsize Time Series Analysis — Final Project\\
Academic Year 2024--2025\par}
\vspace{1.5cm}
{\normalsize\textit{Email: [your email]}\par}
\vfill
{\normalsize\today\par}
\end{titlepage}

% ===========================================================================
\begin{abstract}
\noindent
The rapid growth of India's domestic aviation sector, combined with opaque and
highly volatile pricing mechanisms employed by airlines, creates significant
information asymmetry for consumers and travel planners.  This paper presents a
comprehensive end-to-end machine-learning and time-series framework for
predicting domestic flight fares, applied to a real-world dataset of 10{,}683
booking records spanning March to June 2019.  Starting from raw ticket data, we
design a twelve-stage pipeline that includes data cleaning, correction of a
systematic duration-parsing error that inflated feature values by up to
$10\times$, feature engineering that expands the original 11 columns into 45
predictive features, and a structured multi-model evaluation protocol.

We benchmark five regression algorithms (Linear Regression, Decision Tree,
Random Forest, XGBoost, LightGBM), conduct both random and temporally-ordered
train/test splits, formally assess the applicability of four time-series
paradigms (GARCH, SVAR, SARIMA, SARIMAX) to this data, fit a
SARIMAX ARIMA(0,1,2) model on weekly price aggregates, construct a hybrid
LightGBM+SARIMAX architecture, apply Bayesian hyperparameter optimisation via
Optuna (50 trials per model), and generate route-level fare forecasts for two
major corridors.  The study produces 25 analytical visualisations and 10
structured output reports.

The tuned XGBoost regressor achieves the best cross-sectional generalisation
($R^2 = 0.8400$, MAE $= \text{Rs.}\ 1{,}162$, RMSE $= \text{Rs.}\ 1{,}827$).
Under a rigorous temporal split (train: March--May, test: June 2019), the hybrid
LightGBM+SARIMAX model raises $R^2$ from 0.5351 to 0.7629---a gain of
$+22.78$ percentage points over the temporal baseline.  Feature importance
analysis identifies number of stops (23.38\%), airline tier (14.53\%), and
journey day-of-month (10.81\%) as the three dominant pricing signals.
Route-level forecasting reveals a pronounced full-service vs.\ budget
pricing divide: Jet Airways averages Rs.\,11{,}255 on the Delhi--Cochin
corridor versus Rs.\,6{,}255 for GoAir---a 1.8$\times$ premium.  This work
demonstrates that combining machine learning with time-series trend signals
substantially improves performance under realistic out-of-sample conditions.
\end{abstract}

\vspace{0.5em}
\noindent\textbf{Keywords:} flight fare prediction, machine learning, LightGBM,
XGBoost, SARIMAX, hybrid model, feature engineering, time series, Indian aviation,
Optuna hyperparameter optimisation

\newpage
\tableofcontents
\newpage

% ===========================================================================
\section{Introduction and Project Motivation}
\label{sec:intro}
% ===========================================================================

\subsection{Background}

India's civil aviation sector has undergone remarkable transformation over the
past decade.  According to the Directorate General of Civil Aviation (DGCA),
domestic passenger traffic grew from 81 million in 2015 to over 144 million in
2019~\cite{dgca2019}, making India the third-largest domestic aviation market
globally.  This growth has been driven by the simultaneous emergence of multiple
low-cost carriers (LCCs)---IndiGo, SpiceJet, GoAir, and Air Asia---competing
alongside established full-service airlines such as Jet Airways, Air India, and
Vistara.

The collision of intense competition and sophisticated revenue management
practices has produced a fare environment that is both highly dynamic and
difficult for consumers to navigate.  Airlines employ yield management systems
that adjust seat prices in real time based on remaining seat inventory, booking
lead time, competitive positioning, and seasonal demand patterns~\cite{belobaba2009}.
As a result, the same economy seat on the same flight can differ by several
thousand rupees depending purely on when the ticket is purchased.

\subsection{Problem Statement}

This project addresses the following core question: \textit{Given observable
features of a flight itinerary (airline, route, departure time, number of stops,
and booking date), can we build a predictive model that estimates the fare with
sufficient accuracy to support consumer decision-making and travel budgeting?}

This is framed as a supervised regression problem.  The target variable is the
ticket price in Indian Rupees (INR).  The modelling challenge is threefold:
\begin{enumerate}
    \item \textbf{Non-linearity}: Fare determination involves complex,
    multiplicative interactions between airline identity, route length,
    stop count, and temporal booking signals that linear models cannot capture.

    \item \textbf{Temporal structure}: Flight prices change systematically over
    time (weekly and monthly seasonality), which means naive cross-sectional
    evaluation overstates real-world model performance.  A model that has seen
    ``future'' booking patterns in its training set appears better than it would
    in deployment.

    \item \textbf{Limited time range}: The dataset covers only four months
    (March--June 2019), which constrains the depth of statistical time-series
    modelling possible.  A key question is whether classical time-series methods
    (SARIMA, SARIMAX) add value given this data availability constraint.
\end{enumerate}

\subsection{Academic Context and Project Scope}

This paper constitutes the final project for the \textit{Time Series Analysis}
course at the National Economics University, Vietnam.  The project requirements
specify that both machine learning and time-series techniques must be applied
and evaluated.  Accordingly, the scope of this project is deliberately broader
than a typical ML fare prediction study: we rigorously assess time-series
model suitability, demonstrate that a hybrid architecture outperforms either
approach alone under temporal conditions, and provide a formal justification for
why classical TS models (GARCH, SVAR) are not appropriate for this data type.

\subsection{Paper Contributions}

The main contributions of this paper are:
\begin{enumerate}
    \item \textbf{Duration-parsing bug identification and correction}: We
    identify and fix a systematic error present in widely circulated public
    implementations of this dataset, which inflated flight-duration feature
    values by up to $10\times$ and corrupted downstream model training.

    \item \textbf{Comprehensive 45-feature engineering}: From 11 raw columns
    we engineer 45 predictive features using cyclical time encodings, expert
    airline tier classification, one-hot encoding, and route feature extraction.

    \item \textbf{Rigorous temporal evaluation}: We demonstrate that temporal
    train/test splitting (March--May train, June test) reduces LightGBM $R^2$
    by 29 percentage points compared to random splitting---a substantial
    over-optimism risk that is rarely discussed in published fare prediction studies.

    \item \textbf{Hybrid ML+SARIMAX architecture}: We show that injecting a
    SARIMAX weekly-trend signal as an additional input feature to LightGBM
    recovers most of the temporal performance gap, improving $R^2$ from 0.5351
    to 0.7629.

    \item \textbf{Time-series suitability analysis}: We formally evaluate GARCH,
    SVAR, SARIMA, and SARIMAX for this data type, providing a reusable framework
    for future practitioners facing similar dataset constraints.

    \item \textbf{Route-level fare forecasting}: We generate airline-stratified
    June 2019 fare predictions for two major domestic corridors, quantifying the
    full-service vs.\ budget pricing divide.
\end{enumerate}

\subsection{Paper Organisation}

The remainder of this paper proceeds as follows.
Section~\ref{sec:related} reviews related work.
Section~\ref{sec:data} describes the dataset and exploratory analysis.
Section~\ref{sec:features} presents the feature engineering pipeline.
Section~\ref{sec:method} details the methodology and rationale for each stage.
Section~\ref{sec:metrics} explains the evaluation metrics.
Section~\ref{sec:results} reports all experimental results with figure-by-figure
analysis.
Section~\ref{sec:discussion} discusses key findings, limitations, and practical
implications.
Section~\ref{sec:conclusion} concludes with future work directions.

% ===========================================================================
\section{Related Work}
\label{sec:related}
% ===========================================================================

\subsection{Machine Learning for Airfare Prediction}

The airfare prediction problem has attracted sustained research attention since
Etzioni et al.~\cite{etzioni2003} introduced the \textit{Hamlet} system, which
used historical fare data to advise travellers on the optimal purchase moment.
Their system achieved a 603\% return on investment versus naive buying, establishing
that airfares contain predictable temporal structure that algorithms can exploit.

Grover~\cite{grover2015} applied linear regression, support vector regression,
and random forests to US domestic fare data, reporting that ensemble methods
outperformed linear models with $R^2$ improvements of 0.18--0.24.  Tziridis
et al.~\cite{tziridis2017} compared na\"{i}ve Bayes, random forest, and neural
networks on European airfare data, finding that tree-based models achieved mean
absolute percentage errors as low as 5\%.  Janssen et al.~\cite{janssen2019}
demonstrated that XGBoost matched or exceeded deep learning on structured
tabular airfare data, consistent with the broader empirical finding that
gradient-boosted trees consistently outperform neural networks on tabular
regression tasks~\cite{grinsztajn2022}.

For the Indian aviation context, Raghav and Sinha~\cite{raghav2020} explored
decision trees and random forests on a similar Kaggle-sourced dataset, achieving
$R^2 \approx 0.80$ without temporal evaluation.  Our work extends this baseline
by introducing temporal splitting, a hybrid time-series component, and systematic
hyperparameter optimisation.

\subsection{Time-Series Methods for Pricing}

Box--Jenkins ARIMA and its seasonal extension SARIMA have been applied to
passenger demand forecasting and capacity planning for decades~\cite{grubb2001,
shumway2017, box2015}.  The SARIMAX extension allows inclusion of exogenous
variables such as holidays and fuel prices~\cite{box2015}, which has shown
gains in retail pricing and hotel revenue management.

Grosche et al.~\cite{grosche2007} demonstrated that demand-side time-series
models complement revenue management systems on long-haul routes.  However,
the effectiveness of SARIMA for short series (under 24 months) is well
documented to be limited~\cite{shumway2017}, which directly motivates our
suitability analysis in Section~\ref{sec:method}.

\subsection{Hybrid Architectures}

Zhang~\cite{zhang2003} introduced the idea of combining ARIMA trend
decomposition with neural network modelling of residuals, showing that the
hybrid frequently outperforms either approach alone.  This principle has been
extended to energy price forecasting, stock markets, and supply chain demand.
Our work is among the first to apply this hybrid paradigm explicitly to airfare
prediction with a formal evaluation of the improvement under temporal conditions.

\subsection{Hyperparameter Optimisation}

Early hyperparameter search used exhaustive grid search or random
search~\cite{bergstra2012}.  Bayesian optimisation with Tree-structured Parzen
Estimators (TPE) was shown by Bergstra et al.~\cite{bergstra2011} to find
better solutions in fewer function evaluations.  Akiba et al.~\cite{akiba2019}
introduced Optuna, a \textit{define-by-run} TPE framework, which has become a
standard component of competitive ML pipelines.  We use Optuna with 50 trials
per model in this study.

% ===========================================================================
\section{Dataset Description and Exploratory Analysis}
\label{sec:data}
% ===========================================================================

\subsection{Raw Dataset}

The dataset was sourced from Kaggle's \textit{Flight Fare Prediction}
repository, originally scraped from the Ease My Trip online booking platform
during the booking season spanning March to June 2019.  It contains
\textbf{10{,}683 flight booking records} with 11 columns, described in
Table~\ref{tab:raw_features}.

\begin{table}[h]
\caption{Raw Dataset Features (11 Columns)}
\label{tab:raw_features}
\centering
\begin{tabular}{llp{7cm}}
\toprule
\textbf{Column} & \textbf{Type} & \textbf{Description} \\
\midrule
Airline          & Categorical & Name of the carrier (12 unique airlines) \\
Date\_of\_Journey & String     & Journey date in DD/MM/YYYY format \\
Source           & Categorical & Departure city (5 cities) \\
Destination      & Categorical & Arrival city (6 cities) \\
Route            & String      & Intermediate waypoints (128 unique routes) \\
Dep\_Time        & String      & Scheduled departure time (HH:MM) \\
Arrival\_Time    & String      & Scheduled arrival time (HH:MM, may include +1 day) \\
Duration         & String      & Flight duration encoded as ``Xh Ym'' \\
Total\_Stops     & String      & ``non-stop'', ``1 stop'', ``2 stops'', \ldots \\
Additional\_Info & Categorical & Ancillary info: meal, baggage, in-flight entertainment \\
Price            & Integer     & Fare in Indian Rupees -- \textbf{prediction target} \\
\bottomrule
\end{tabular}
\end{table}

The dataset covers 12 airlines, 5 source cities
(Delhi, Mumbai, Bangalore, Kolkata, Chennai), 6 destination cities
(Cochin, Delhi, Hyderabad, Kolkata, Mumbai, New Delhi), and 128 unique route
combinations.  The booking dates span 44 unique journey dates across the
March--June 2019 window.

\subsection{Target Variable: Price}

The \textit{Price} variable ranges from Rs.\,1{,}759 (ultra-budget connecting
flights) to Rs.\,79{,}512 (long-haul full-service fares), with summary
statistics:
\begin{itemize}
    \item \textbf{Mean}: Rs.\,9{,}087
    \item \textbf{Median}: Rs.\,8{,}372
    \item \textbf{Standard Deviation}: Rs.\,4{,}611
    \item \textbf{25th percentile}: Rs.\,5{,}277
    \item \textbf{75th percentile}: Rs.\,12{,}373
\end{itemize}

The median being lower than the mean (Rs.\,8{,}372 vs.\ Rs.\,9{,}087) confirms
a right-skewed distribution.  This skew is driven by a small number of very
expensive tickets---primarily Jet Airways and Air India fares on longer
connecting routes.  The interquartile range of Rs.\,7{,}096 is large relative
to the mean, indicating high heterogeneity in pricing even within the same
route and airline class.

\subsection{Figure 01: Price Distribution}

The price distribution histogram (Fig.\,01) reveals a unimodal, right-skewed
distribution with the mode near Rs.\,5{,}000--6{,}000.  The long right tail
(fares above Rs.\,40{,}000) represents fewer than 2\% of bookings but pulls
the mean substantially above the median.  This skew has two practical
implications: (1) models trained with squared-error loss may over-prioritise
accurate prediction of expensive tickets, and (2) RMSE will be more sensitive
to outliers than MAE, making the MAE/RMSE comparison informative about the
magnitude of large prediction errors (see Section~\ref{sec:metrics}).

\begin{figure}[htbp]
\centering
\includegraphics[width=0.75\textwidth]{01_price_distribution.png}
\caption{Figure 01 --- Distribution of flight fares (INR). Right-skewed with mode near Rs.\,5{,}000--6{,}000 and long right tail to Rs.\,79{,}512.}
\label{fig:f01}
\end{figure}

\subsection{Figure 02: Categorical Variable Distributions}

Count plots (Fig.\,02) show the frequency of each level within categorical
features:
\begin{itemize}
    \item \textbf{Airline}: Jet Airways leads with 3{,}849 records (36.1\%),
    reflecting its position as India's dominant full-service carrier before
    its April 2019 shutdown.  IndiGo is second with approximately 2{,}100
    records (20\%).  Smaller carriers (Air Asia, TruJet) have fewer than
    500 records each.
    \item \textbf{Source}: Delhi accounts for 4{,}537 records (42\%),
    followed by Mumbai and Bangalore.  This skew means predictions for
    Delhi departures will be more reliable than for Chennai.
    \item \textbf{Destination}: Cochin and Delhi are the two dominant
    destinations, together accounting for over 85\% of records.
    \item \textbf{Total\_Stops}: 1-stop itineraries dominate (5{,}625
    records, 53\%), followed by non-stop (3{,}212 records, 30\%).
    Two-stop and above together represent only 17\% of bookings.
\end{itemize}

\begin{figure}[htbp]
\centering
\includegraphics[width=\textwidth]{02_categorical_distributions.png}
\caption{Figure 02 --- Count distributions for categorical features. Jet Airways and Delhi-origin bookings dominate the dataset; 1-stop itineraries account for 53\% of records.}
\label{fig:f02}
\end{figure}

\subsection{Figure 03: Categorical Variables vs.\ Price}

Box plots of price stratified by each categorical variable (Fig.\,03) reveal
the pricing hierarchy directly:
\begin{itemize}
    \item \textbf{By Airline}: Jet Airways and Air India show medians near
    Rs.\,10{,}800 and Rs.\,9{,}500 respectively, approximately
    $2\times$--$2.5\times$ higher than budget carriers IndiGo (Rs.\,4{,}600)
    and GoAir (Rs.\,4{,}200).  Vistara occupies a middle position
    (Rs.\,6{,}500--7{,}000) despite being a full-service airline, likely
    due to its entry-level pricing strategy to capture market share.
    \item \textbf{By Source}: Delhi departures command a small premium over
    other sources, consistent with higher demand on the hub.
    \item \textbf{By Destination}: Routes to Cochin are on average more
    expensive than short-haul routes to Delhi, reflecting route length effects.
\end{itemize}

\begin{figure}[htbp]
\centering
\includegraphics[width=\textwidth]{03_categorical_vs_price.png}
\caption{Figure 03 --- Box plots of price by airline, source, and destination. Full-service carriers (Jet Airways, Air India) show medians 2--2.5$\times$ higher than budget LCCs.}
\label{fig:f03}
\end{figure}

\subsection{Figure 04: Total Stops vs.\ Price}

The stops-versus-price relationship (Fig.\,04) is somewhat counterintuitive:
non-stop flights are \textit{not} systematically cheaper than 1-stop
itineraries.  In fact, direct flights from Delhi to Cochin (a 3-hour
non-stop) are priced higher than 1-stop alternatives because the non-stop
service is a premium product commanded by full-service carriers.  However,
at 2+ stops, prices decline as passengers trade time for cost.  This
non-monotonic pattern is a key reason linear models underperform---the
relationship between stops and price depends on the airline identity
(a non-stop IndiGo flight is cheaper than a 1-stop Jet Airways flight).

\begin{figure}[htbp]
\centering
\includegraphics[width=0.7\textwidth]{04_stops_vs_price.png}
\caption{Figure 04 --- Price distribution by number of stops. Non-stop flights are not necessarily cheaper due to the full-service airline premium on direct long-haul routes.}
\label{fig:f04}
\end{figure}

\subsection{Figure 05: Numerical Variable Distributions}

Histograms of numerical features (Fig.\,05) show:
\begin{itemize}
    \item \textbf{Departure Hour}: Bi-modal distribution with peaks at
    06:00--08:00 (morning business travel) and 18:00--20:00 (evening
    commuter traffic).  Off-peak hours (noon, midnight) have lower flight
    density, suggesting less competition and potentially different pricing.
    \item \textbf{Departure Minute}: Multi-modal at 0, 15, 30, 45 --- consistent
    with airline scheduling conventions.
    \item \textbf{Duration}: Majority of domestic flights fall in the
    2--5 hour range.  A small cluster exists at 15--47 hours, corresponding
    to heavily connected long-haul domestic itineraries.
\end{itemize}

\begin{figure}[htbp]
\centering
\includegraphics[width=\textwidth]{05_numerical_distributions.png}
\caption{Figure 05 --- Histograms of numerical features. Departure hours show a bi-modal pattern (morning and evening peaks); flight duration clusters in the 2--5 hour domestic range.}
\label{fig:f05}
\end{figure}

\subsection{Figure 06: Correlation Heatmap}

The correlation heatmap (Fig.\,06) of numerical features after preprocessing
reveals several important relationships:
\begin{itemize}
    \item \textbf{Total\_Stops vs.\ Price}: Moderate positive correlation
    ($\rho \approx 0.22$).  More stops correlate with slightly higher price
    because stop-count covaries with route length, and longer routes cost more
    even with connections.
    \item \textbf{Duration\_Hours vs.\ Price}: Positive correlation
    ($\rho \approx 0.35$), the strongest pair in the matrix.  Longer flights
    cost more, driven by route economics.
    \item \textbf{Departure timing features}: Low correlation with price
    individually, but these features interact non-linearly with airline and
    route (morning Jet Airways Delhi--Mumbai fares are not simply the sum of
    a ``morning premium'' and a ``Jet Airways premium'').
\end{itemize}

The heatmap thus motivates the use of non-linear models (tree ensembles) over
linear regression, which can only model additive main effects.

\begin{figure}[htbp]
\centering
\includegraphics[width=0.8\textwidth]{06_correlation_heatmap.png}
\caption{Figure 06 --- Correlation heatmap of numerical features. Duration\_Hours shows the strongest linear association with Price ($\rho \approx 0.35$); most features interact non-linearly.}
\label{fig:f06}
\end{figure}

\subsection{Figure 07: Monthly Price Trend}

The monthly mean price trend (Fig.\,07) shows:
\begin{itemize}
    \item March: mean Rs.\,8{,}700
    \item April: mean Rs.\,8{,}900 (slight uplift; possible Jet Airways
    pricing anomalies pre-shutdown)
    \item May: mean Rs.\,9{,}100 (summer travel season beginning)
    \item June: mean Rs.\,9{,}610 (early-monsoon peak)
\end{itemize}

The upward trend from March to June (approximately +10\% over the period)
motivates the time-series component of this project.  If a model is trained
on March--May patterns and asked to predict June prices, it will tend to
underestimate unless it has access to the upward trend signal.  This is
precisely the gap that the hybrid SARIMAX component addresses.

\begin{figure}[htbp]
\centering
\includegraphics[width=0.8\textwidth]{07_monthly_price_trend.png}
\caption{Figure 07 --- Monthly mean fare trend (March--June 2019). Prices rise approximately 10\% from March (Rs.\,8{,}700) to June (Rs.\,9{,}610), motivating the temporal split and SARIMAX component.}
\label{fig:f07}
\end{figure}

% ===========================================================================
\section{Data Preprocessing}
\label{sec:preproc}
% ===========================================================================

\subsection{Missing Value Handling}

Of the 10{,}683 raw records, two contained missing values: one in the
\textit{Route} column and one in \textit{Total\_Stops}.  Because the
missingness was isolated (0.02\%) and non-informative (neither row belonged to
a particularly rare airline or route combination), both were removed rather
than imputed.  This left \textbf{10{,}462 clean records} for all downstream
processing.

\subsection{City Name Standardisation}

The \textit{Destination} column contained two representations for the same
city: ``New Delhi'' and ``Delhi''.  All instances of ``New Delhi'' were
replaced with ``Delhi'' to avoid the model treating these as two distinct
destinations.  Failing to perform this standardisation would create a
spurious extra feature column in one-hot encoding and split route-level
signal across two dummy variables.

\subsection{Duration Parsing: Identification and Correction of a Systematic Error}
\label{sec:bug}

The \textit{Duration} column encodes flight duration as a formatted string,
for example ``2h 50m'' (2 hours, 50 minutes).  This string must be parsed into
a numeric feature representing total hours.  A widely circulated implementation
of this parsing used the following Python code:

\begin{verbatim}
train_data['Duration_hours'] = train_data['Duration']
    .str.replace('h', '1')
    .str.replace('m', '')
    .apply(eval)
\end{verbatim}

This code contains a critical logical error.  When the character `h' is
replaced by the character `1', the string ``2h 50m'' becomes ``21 50'', and
calling \texttt{eval()} on this string evaluates it as the Python arithmetic
expression $21 \times 50 = 1050$.  Table~\ref{tab:bugfix} illustrates the
impact for several representative values.

\begin{table}[h]
\caption{Impact of Duration Parsing Error}
\label{tab:bugfix}
\centering
\begin{tabular}{lrr}
\toprule
\textbf{Raw Duration} & \textbf{Buggy Value (hours)} & \textbf{Correct Value (hours)} \\
\midrule
2h 50m  & 1050  & 2.83 \\
1h 30m  & 30    & 1.50 \\
5h 00m  & 500   & 5.00 \\
12h 45m & 1245  & 12.75 \\
\bottomrule
\end{tabular}
\end{table}

The errors span a factor of $10\times$--$400\times$.  Because flight duration
has the 5th-highest predictive importance in our model (8.48\%), this bug would
severely corrupt any model trained using the erroneous values.

The corrected implementation uses regular expressions to independently extract
hours and minutes:

\begin{verbatim}
import re

def _parse_duration(s: str) -> float:
    h = re.search(r'(\d+)h', s)
    m = re.search(r'(\d+)m', s)
    return (int(h.group(1)) if h else 0) + \
           (int(m.group(1)) if m else 0) / 60.0
\end{verbatim}

This yields \texttt{Duration\_Total\_Hour} values in the range $[0.25, 47.17]$
hours, consistent with the actual domestic flight schedule.  Short-duration
values (0.25h $\approx$ 15 minutes) correspond to brief shuttle-style legs
within major route chains.

% ===========================================================================
\section{Feature Engineering}
\label{sec:features}
% ===========================================================================

Feature engineering is the process of transforming raw input columns into
representations that are more informative, more separable, and more compatible
with the modelling algorithm.  Starting from 11 raw columns and 13 cleaned
attributes, the engineering pipeline produces \textbf{45 predictive features}.

\subsection{Date and Time Decomposition}

The raw \textit{Date\_of\_Journey} string (``DD/MM/YYYY'') is parsed into:
\begin{itemize}
    \item \textbf{Journey\_Day} (1--31): Day of the month.  With importance of
    10.81\%, this is the third most predictive feature, capturing the fact that
    airlines reprice inventory aggressively as the departure date approaches
    within the month.
    \item \textbf{Journey\_Month} (3--6): The booking month.  Importance of
    1.93\%, lower because the dataset only spans four months with limited
    inter-month price variance compared to annual seasonality.
    \item \textbf{Journey\_Weekdays} (0=Monday--6=Sunday): Day of the week.
    Importance of 1.33\%; Friday and Sunday departures tend to be slightly
    more expensive due to business and leisure travel peaks.
    \item \textbf{Is\_Weekend} (binary): Boolean flag for Saturday/Sunday.
    Importance of 0.68\%.
\end{itemize}

Similarly, \textit{Dep\_Time} and \textit{Arrival\_Time} are parsed into
\textbf{Dep\_Hour}, \textbf{Dep\_Min}, \textbf{Arrival\_Hours}, and
\textbf{Arrival\_Min}.

\subsection{Cyclical Time Encoding}

Ordinal encoding of cyclic variables (month, weekday, hour) is problematic:
a model would treat January (month=1) and December (month=12) as maximally
different, when they are in fact adjacent in the annual calendar.  Cyclical
encoding resolves this by projecting the ordinal value onto a unit circle:
\begin{align}
    \text{Feature\_Sin} &= \sin\!\left(\frac{2\pi k}{P}\right) \\
    \text{Feature\_Cos} &= \cos\!\left(\frac{2\pi k}{P}\right)
\end{align}
where $k$ is the ordinal value and $P$ is the period
(12 for months, 7 for weekdays, 24 for hours).

This produces four pairs of features:
\texttt{Month\_Sin/Cos}, \texttt{Weekday\_Sin/Cos},
\texttt{Dep\_Hour\_Sin/Cos}, and \texttt{Arr\_Hour\_Sin/Cos}.
The cosine component of departure hour (\texttt{Dep\_Hour\_Cos}) ranks 10th
in importance (1.64\%), confirming that cyclical hour information adds signal
beyond the raw numeric hour.

\subsection{Departure Period Categorisation}

The 24-hour clock is binned into five named periods:
\begin{itemize}
    \item \textbf{Pre-Dawn} ($< 06{:}00$): very early red-eye departures
    \item \textbf{Morning} ($06{:}00$--$11{:}59$): peak business travel window
    \item \textbf{Afternoon} ($12{:}00$--$16{:}59$): mid-day travel
    \item \textbf{Evening} ($17{:}00$--$20{:}59$): peak commuter window
    \item \textbf{Night} ($\geq 21{:}00$): late-night departures
\end{itemize}

Both departure period (\texttt{Dep\_Period}) and arrival period
(\texttt{Arr\_Period}) are created.  These allow tree-based models to learn
non-linear pricing patterns at the session level rather than needing to
identify period boundaries from raw minute values.

\subsection{Airline Tier Classification}

Each airline is classified into one of two tiers based on service level:
\begin{itemize}
    \item \textbf{Full-Service (tier = 1)}: Jet Airways, Air India, Vistara.
    These airlines provide complimentary meals, checked baggage, and premium
    cabin options.
    \item \textbf{Budget / LCC (tier = 0)}: IndiGo, SpiceJet, GoAir, Air Asia,
    TruJet, and multi-carrier codeshare combinations.
    These carriers use unbundled pricing models.
\end{itemize}

\texttt{Airline\_Tier} achieves 14.53\% importance---the second-highest feature
in the model.  This is remarkable because it is a binary variable derived
purely from expert domain knowledge, yet it outperforms all individual airline
one-hot dummies (the Jet Airways dummy ranks 4th at 9.29\%).  The tier feature
compresses within-class variance and captures the premium vs.\ budget pricing
structure at a higher abstraction level.

\subsection{Stop Count Ordinal Encoding}

The \textit{Total\_Stops} string column (``non-stop'', ``1 stop'', ``2 stops'',
``3 stops'', ``4 stops'') is mapped to the ordinal integer representation
(0, 1, 2, 3, 4).  This is the most important feature overall at 23.38\%.

\subsection{One-Hot Encoding}

Three categorical columns are one-hot encoded:
\begin{itemize}
    \item \textbf{Airline}: 11 binary columns (Jet Airways, IndiGo, Air India,
    SpiceJet, GoAir, Vistara, Air Asia, Multiple carriers, TruJet,
    Vistara, and one reference category absorbed).
    \item \textbf{Source}: 4 binary columns (one reference city dropped).
    \item \textbf{Destination}: 5 binary columns (one reference city dropped).
\end{itemize}

\textit{Route} features are also extracted: two route-level binary columns
(\texttt{Route\_Banglore\_to\_Delhi}, \texttt{Route\_Mumbai\_to\_Hyderabad},
\texttt{Route\_Delhi\_to\_Cochin}, etc.) capture major corridor effects beyond
what source and destination alone encode.

\subsection{Complete Feature Set Summary}

The complete 45-feature set with importance ranks is shown in
Table~\ref{tab:all_features}.  The top 11 features collectively explain 79.5\%
of model split decisions, and the remaining 34 features contribute the
remaining 20.5\%, capturing fine-grained route and airline effects.

\begin{table}[h]
\caption{All 45 Features by LightGBM Importance (full list)}
\label{tab:all_features}
\centering
\small
\begin{tabular}{rlrr}
\toprule
\textbf{Rank} & \textbf{Feature} & \textbf{Importance (\%)} & \textbf{Cumulative (\%)} \\
\midrule
1  & Total\_Stops               & 23.38 & 23.38 \\
2  & Airline\_Tier              & 14.53 & 37.91 \\
3  & Journey\_Day               & 10.81 & 48.72 \\
4  & Airline\_Jet Airways       &  9.29 & 58.01 \\
5  & Duration\_Hours            &  8.48 & 66.49 \\
6  & Dep\_Min                   &  2.72 & 69.21 \\
7  & Weekday\_Cos               &  2.63 & 71.83 \\
8  & Month\_Cos                 &  2.49 & 74.33 \\
9  & Journey\_Month             &  1.93 & 76.25 \\
10 & Dep\_Hour\_Cos             &  1.64 & 77.90 \\
11 & Arrival\_Min               &  1.60 & 79.50 \\
12 & Arrival\_Hours             &  1.57 & 81.07 \\
13 & Dep\_Hours                 &  1.53 & 82.59 \\
14 & Journey\_Weekdays          &  1.33 & 83.93 \\
15 & Arr\_Hour\_Cos             &  1.32 & 85.25 \\
16 & Dep\_Hour\_Sin             &  1.25 & 86.49 \\
17 & Month\_Sin                 &  1.10 & 87.59 \\
18 & Weekday\_Sin               &  1.09 & 88.68 \\
19 & Source\_Banglore           &  1.08 & 89.76 \\
20 & Arr\_Hour\_Sin             &  1.01 & 90.77 \\
21 & Dep\_Period                &  0.98 & 91.76 \\
22 & Arr\_Period                &  0.93 & 92.69 \\
23 & Airline\_Multiple carriers &  0.83 & 93.51 \\
24 & Route\_Banglore\_to\_Delhi &  0.74 & 94.26 \\
25 & Is\_Weekend                &  0.68 & 94.94 \\
26 & Destination\_Delhi         &  0.68 & 95.62 \\
27 & Destination\_Hyderabad     &  0.59 & 96.21 \\
28 & Route\_Mumbai\_to\_Hyd.   &  0.49 & 96.70 \\
29 & Source\_Mumbai             &  0.47 & 97.17 \\
30 & Airline\_IndiGo            &  0.40 & 97.58 \\
31--44 & [Remaining features]  & 2.42  & 100.00 \\
\bottomrule
\end{tabular}
\end{table}

% ===========================================================================
\section{Evaluation Metrics}
\label{sec:metrics}
% ===========================================================================

This section defines and motivates the four evaluation metrics used throughout
the study.  Understanding the properties of each metric is essential for
interpreting the model comparison tables correctly.

\subsection{Coefficient of Determination ($R^2$)}

\begin{equation}
R^2 = 1 - \frac{\text{SS}_{\text{res}}}{\text{SS}_{\text{tot}}} = 1 - \frac{\sum_{i=1}^{n}(y_i - \hat{y}_i)^2}{\sum_{i=1}^{n}(y_i - \bar{y})^2}
\end{equation}

$R^2$ measures the proportion of variance in the target variable (price)
that is explained by the model.  It is scale-free, ranging from $(-\infty, 1]$,
where:
\begin{itemize}
    \item $R^2 = 1$: perfect prediction
    \item $R^2 = 0$: the model performs as well as simply predicting the mean
    fare for every observation
    \item $R^2 < 0$: the model is \textit{worse} than predicting the mean (this
    occurs in the SARIMAX weekly forecast: $R^2 = -0.3701$)
\end{itemize}

\textbf{Advantage}: Intuitive and universally comparable across datasets.\\
\textbf{Limitation}: $R^2$ always increases (or stays the same) when more
features are added, even if they add no real information.  This necessitates
adjusted $R^2$.

\subsection{Adjusted $R^2$}

\begin{equation}
\bar{R}^2 = 1 - (1 - R^2)\frac{n-1}{n-p-1}
\end{equation}

where $n$ is the number of test observations and $p$ is the number of features.
Adjusted $R^2$ penalises model complexity, decreasing when features that do
not improve predictive power are added.  With 45 features and $n \approx 2{,}093$
test observations, the penalty is small (approximately 0.008), so $R^2$ and
$\bar{R}^2$ are close in this study.  For the best model (LightGBM baseline),
$R^2 = 0.8275$ vs.\ $\bar{R}^2 = 0.8237$.

\subsection{Mean Absolute Error (MAE)}

\begin{equation}
\text{MAE} = \frac{1}{n}\sum_{i=1}^{n}|y_i - \hat{y}_i|
\end{equation}

MAE reports the average magnitude of prediction error in the original unit
(Indian Rupees).  It is robust to outliers---a single ticket priced at
Rs.\,79{,}512 that is predicted at Rs.\,50{,}000 contributes only
Rs.\,29{,}512 to the MAE sum, the same as two ordinary tickets each wrong
by Rs.\,14{,}756.

\textbf{Practical interpretation}: A MAE of Rs.\,1{,}162 for the tuned XGBoost
model means that, on average, the model's fare estimate is off by Rs.\,1{,}162.
At the 2019 RBI exchange rate of Rs.\,70.42/USD, this is \textbf{approximately
USD 16.50}---a meaningful but tolerable error for personal travel budgeting.

\subsection{Root Mean Squared Error (RMSE)}

\begin{equation}
\text{RMSE} = \sqrt{\frac{1}{n}\sum_{i=1}^{n}(y_i - \hat{y}_i)^2}
\end{equation}

RMSE penalises large errors quadratically: an error of Rs.\,5{,}000 counts
$25\times$ more than an error of Rs.\,1{,}000 ($5000^2 / 1000^2 = 25$).
This makes RMSE more sensitive to outlier predictions.

\textbf{Key diagnostic: RMSE/MAE ratio}. When RMSE $\gg$ MAE, a few large
errors dominate the RMSE.  For our best models:
\begin{itemize}
    \item Tuned XGBoost: RMSE/MAE $= 1827/1162 = 1.57$
    \item Tuned LightGBM: RMSE/MAE $= 1843/1182 = 1.56$
    \item Tuned Random Forest: RMSE/MAE $= 1944/1233 = 1.58$
\end{itemize}

All three models show ratios of approximately 1.57.  A ratio above 1.5 indicates
the presence of some large prediction errors---likely corresponding to premium
Jet Airways tickets on unusual routes where training data is sparse.  This ratio
is consistent across models, suggesting the large errors are data-driven
(inherently hard cases) rather than model-specific failures.

% ===========================================================================
\section{Methodology}
\label{sec:method}
% ===========================================================================

The pipeline consists of twelve sequential stages.  Each stage is described
below with its purpose, implementation, and the reasoning behind the design choice.

\subsection{Stage 1: Random 80/20 Train/Test Split}

\textbf{Purpose}: Establish a cross-sectional performance baseline.

The 10{,}462 cleaned records are split randomly into:
\begin{itemize}
    \item \textbf{Training set}: 8{,}369 records (80\%)
    \item \textbf{Test set}: 2{,}093 records (20\%)
\end{itemize}

\textbf{Rationale}: Random splitting is the standard evaluation protocol for
cross-sectional (non-temporal) prediction tasks.  It ensures that the test
set is drawn from the same distribution as training, providing an estimate
of how well the model would generalise to randomly selected future bookings.

\textbf{Limitation}: Random splitting ignores temporal structure.  If a
booking from June 3 appears in the training set alongside a booking from
March 15, the model has implicitly ``seen'' June pricing patterns during
training.  This inflates estimated performance compared to real-world
deployment, where predictions are always made about future, unseen dates.

The split results are serialised to
\texttt{outputs/processed/X\_train.csv}, \texttt{X\_test.csv},
\texttt{y\_train.csv}, and \texttt{y\_test.csv} for reproducibility.

\subsection{Stage 2: Feature Importance Analysis (Figures 08--09)}
\label{sec:fi}

\textbf{Purpose}: Identify which of the 45 features drive model predictions.

A LightGBM model is trained on the full training set using default parameters.
Feature importance is extracted using the \texttt{split} metric: the number of
times each feature is used as a decision node across all trees, normalised
to sum to 100\%.

\textbf{Why LightGBM for importance?}  LightGBM's leaf-wise tree growth
tends to produce more concentrated feature importance scores than Random
Forest (which spreads importance more uniformly).  This makes the top-feature
hierarchy more interpretable.

\textbf{Figure 08: Feature Importance Bar Chart} shows the individual
importance of each of the 45 features, sorted in descending order.  The
dramatic drop-off after the first five features (from 23.38\% for
\textit{Total\_Stops} to 2.72\% for \textit{Dep\_Min}) indicates that a
small number of features carry most of the predictive signal.

\textbf{Figure 09: Cumulative Feature Importance} plots the cumulative sum of
importance as features are added.  The 80\% cumulative threshold is crossed
at approximately the 11th feature (\textit{Arrival\_Min}, cumulative 79.50\%).
This means that 11 of the 45 features account for four-fifths of the model's
decision-making---the remaining 34 features collectively contribute only 20.5\%.
This insight could be used to build a simplified model for low-latency
prediction contexts.

\begin{figure}[htbp]
\centering
\includegraphics[width=\textwidth]{08_feature_importance.png}
\caption{Figure 08 --- LightGBM split-based feature importance (45 features). Total\_Stops (23.38\%) and Airline\_Tier (14.53\%) are far more important than any other individual feature.}
\label{fig:f08}
\end{figure}

\begin{figure}[htbp]
\centering
\includegraphics[width=0.75\textwidth]{09_cumulative_feature_importance.png}
\caption{Figure 09 --- Cumulative feature importance. The 80\% threshold is crossed at the 11th feature, meaning 11 of 45 features drive four-fifths of all model decisions.}
\label{fig:f09}
\end{figure}

\subsection{Stage 3: Baseline Model Comparison (Figure 10)}

\textbf{Purpose}: Establish performance benchmarks for five regression families.

Five regression algorithms are trained on the 80/20 random split and evaluated
on the 20\% hold-out set.

\textbf{Why these five models?}  They span the complexity spectrum from
linear (no interactions) to deep ensembles:
\begin{enumerate}
    \item \textbf{Linear Regression} (OLS): The null hypothesis of the
    modelling exercise.  If it performs well, the relationships are largely
    additive and linear.  Its poor performance ($R^2 = 0.6037$) confirms
    that non-linearity is essential.

    \item \textbf{Decision Tree} (CART, unlimited depth): A non-linear
    baseline that can capture interactions but is prone to overfitting.
    The extreme train/test gap (0.9697 vs.\ 0.6277) confirms severe
    overfitting.

    \item \textbf{Random Forest}: Reduces Decision Tree variance via
    bagging (bootstrap aggregating) and feature sub-sampling.  Strong
    generalisation ($R^2 = 0.8052$) with controlled overfitting.

    \item \textbf{XGBoost}: Gradient-boosted trees with regularisation
    (L1/L2 penalties, subsampling).  $R^2 = 0.7908$---slightly below
    Random Forest under default settings, but improves substantially
    after Optuna tuning.

    \item \textbf{LightGBM}: Microsoft's leaf-wise gradient boosting
    implementation, which is faster and often achieves better RMSE
    than XGBoost on right-skewed targets.  Best baseline performance:
    $R^2 = 0.8275$, RMSE = Rs.\,1{,}896.
\end{enumerate}

\textbf{Figure 10: Baseline Model Comparison} presents three grouped bar
charts (one each for $R^2$, MAE, and RMSE), sorted by best-to-worst
performance.  LightGBM occupies the top position in $R^2$ and RMSE,
while Random Forest has a slightly lower MAE (Rs.\,1{,}187 vs.\
Rs.\,1{,}261).  The Decision Tree's tall bar in both MAE and RMSE
visualises the cost of overfitting without pruning.

\begin{figure}[htbp]
\centering
\includegraphics[width=\textwidth]{10_baseline_model_comparison.png}
\caption{Figure 10 --- Baseline model comparison across $R^2$, MAE, and RMSE. LightGBM leads in $R^2$ and RMSE; Random Forest achieves the lowest MAE among default configurations.}
\label{fig:f10}
\end{figure}

\subsection{Stage 4: Time-Series Suitability Analysis}

\textbf{Purpose}: Formally evaluate whether classical TS models are appropriate
for this dataset before applying them.

This stage is motivated by a common error in applied machine learning:
applying a modelling technique without verifying whether its assumptions
are satisfied.  We evaluate four paradigms:

\subsubsection{GARCH}

GARCH (Generalised Autoregressive Conditional Heteroskedasticity) models
\textit{volatility clustering} in financial log-returns: periods of high
variance are followed by periods of high variance.  It operates on series
where the variance of the error is time-varying and autocorrelated.

Flight fares are \textit{price levels}, not log-returns.  They do not exhibit
the variance-clustering structure GARCH requires.  Applying GARCH to price
levels would produce meaningless results.  \textbf{Verdict: not applicable.}

\subsubsection{SVAR}

Structural Vector Autoregression (SVAR) models \textit{multiple co-integrated
time series} with structural causal relationships (e.g., fuel price
simultaneously caused by and causing seat demand).  It requires at least two
interrelated time series and a theoretical model of their causal structure.

Our dataset has a single price target and no co-integrated companion series.
Moreover, 119 days of data is far too short for reliable SVAR estimation,
which typically requires 5+ years of data to identify structural impulse
response functions.  \textbf{Verdict: not applicable.}

\subsubsection{SARIMA}

SARIMA requires sufficient data to estimate seasonal parameters reliably.
The Augmented Dickey-Fuller (ADF) test on the daily price series yields
statistic $-4.621$ ($p = 0.000118$, critical value at 5\%: $-2.886$),
confirming the series is stationary.  However:
\begin{itemize}
    \item Annual seasonality ($s = 365$) requires $\geq 2$ full cycles
    ($\geq 730$ days).  Our window of 119 days covers only 16\%.
    \item Even weekly seasonality ($s = 7$) with 17 complete cycles is
    marginal for parameter estimation.
    \item Critically, SARIMA must aggregate individual flights to a daily
    mean price, discarding the airline, route, and stop-count features
    that drive 80\%+ of price variance.
\end{itemize}
\textbf{Verdict: limited standalone value; used only as trend signal.}

\subsubsection{SARIMAX}

SARIMAX adds exogenous regressors to SARIMA.  Month dummies or holiday flags
can be included, providing marginal improvement over pure SARIMA.  However,
the 4-month window constraint remains binding, and the individual-flight
features that tree models exploit are still unavailable.  \textbf{Verdict:
marginal improvement; used in hybrid role as a weekly trend smoother.}

\subsection{Stage 5: Temporal Train/Test Split}

\textbf{Purpose}: Evaluate model performance under realistic deployment
conditions.

The chronological split is:
\begin{itemize}
    \item \textbf{Training set}: March 1 -- May 31, 2019 ($n = 7{,}151$ flights, 68.4\%)
    \item \textbf{Test set}: June 1 -- June 27, 2019 ($n = 3{,}311$ flights, 31.6\%)
\end{itemize}

\textbf{Rationale}: In real deployment, a booking prediction system is always
predicting fares on \textit{future} dates using a model trained on \textit{past}
data.  The temporal split mimics this condition: the model sees March--May
pricing patterns and must generalise to June.  The performance drop from the
random split reveals the magnitude of distributional shift between spring and
early-summer booking patterns.

LightGBM is used as the ML model for this stage because its compiled C++
backend allows training in seconds, enabling rapid iteration over the temporal
split methodology.

\subsection{Stage 6: SARIMAX on Weekly Price Aggregates (Figures 12--17)}
\label{sec:sarimax_method}

\textbf{Purpose}: Fit a time-series model on the coarsest temporal granularity
where it is statistically feasible.

Individual-flight prices cannot be modelled as a time series because multiple
flights with very different prices depart on the same day.  We therefore
aggregate to \textbf{weekly mean prices}, producing 18 weekly observations
spanning February 25 -- June 24, 2019.

\textbf{Figure 12: Daily Price Trend} shows the raw daily mean price as a
time series over the 119-day window.  The trend is noisy at the daily level
(caused by variation in which airlines and routes happen to depart on a
given day) but shows a gradual upward drift from March to June.  There is
no clear weekly periodicity visible in the daily data.

\textbf{Figure 13: Seasonal Decomposition} applies STL (Seasonal and Trend
decomposition using Loess) to the daily series.  The decomposition extracts:
\begin{itemize}
    \item \textbf{Trend component}: Confirms the March--June upward drift
    ($\approx +$Rs.\,1{,}000 over the period).
    \item \textbf{Seasonal component}: Shows a weak weekly pattern; Fridays
    and Sundays have slightly elevated fares.
    \item \textbf{Residual component}: Large residuals indicate that most
    day-to-day price variation is idiosyncratic (driven by which routes and
    airlines are sampled), not seasonal.
\end{itemize}
The large residuals confirm that aggregated daily prices are a poor basis for
time-series modelling---further motivating the weekly aggregation level.

\begin{figure}[htbp]
\centering
\includegraphics[width=0.8\textwidth]{12_daily_price_trend.png}
\caption{Figure 12 --- Daily mean price time series (March--June 2019). Noisy day-to-day variation obscures the underlying upward trend; weekly aggregation is used for SARIMAX.}
\label{fig:f12}
\end{figure}

\begin{figure}[htbp]
\centering
\includegraphics[width=\textwidth]{13_seasonal_decomposition.png}
\caption{Figure 13 --- STL seasonal decomposition of the daily price series. Trend component confirms the March--June uplift; large residuals indicate most variance is cross-sectional, not temporal.}
\label{fig:f13}
\end{figure}

\textbf{Figure 14: ACF/PACF of Daily Series} plots the Autocorrelation Function
(ACF) and Partial Autocorrelation Function (PACF) for the daily price series.
ACF shows significant autocorrelation at lag 1 (consecutive days are correlated)
and some significance at lag 7 (weekly pattern), but most higher-order lags are
not significant.  The PACF cuts off after lag 1, suggesting an AR(1) structure
is dominant at the daily level.  These plots motivated the choice to use
ARIMA rather than purely seasonal models.

\begin{figure}[htbp]
\centering
\includegraphics[width=\textwidth]{14_acf_pacf.png}
\caption{Figure 14 --- ACF and PACF of the daily price series. Significant autocorrelation at lag 1 and lag 7; PACF cuts off after lag 1, suggesting AR(1) structure dominates at the daily level.}
\label{fig:f14}
\end{figure}

\textbf{Auto-ARIMA model selection}: Using the \texttt{pmdarima} library with
AIC criterion, $d \in \{0,1\}$, and $p,q \leq 3$, the algorithm selected
\textbf{ARIMA(0,1,2)} on the 18 weekly observations.

\textbf{Figure 16: ACF/PACF of Weekly Series} plots the autocorrelation of
the 18-observation weekly price series.  ACF shows no significant spikes
beyond lag 1, while PACF shows a marginal spike at lag 2.  The ARIMA(0,1,2)
selection (an IMA model) is consistent with: differencing the series once
removes the stochastic trend, and two moving-average terms capture the remaining
short-memory structure.  The MA(2) coefficients absorb the residual
autocorrelation that a pure random-walk model would miss.

\begin{figure}[htbp]
\centering
\includegraphics[width=\textwidth]{16_weekly_acf_pacf.png}
\caption{Figure 16 --- ACF and PACF of the 18-observation weekly aggregated price series. No significant spikes beyond lag 1--2; consistent with the ARIMA(0,1,2) IMA model selected by auto-ARIMA.}
\label{fig:f16}
\end{figure}

The fitted ARIMA(0,1,2) model on the differenced weekly series is:
\begin{equation}
(1 - B)\,P_t = \theta_1\,\varepsilon_{t-1} + \theta_2\,\varepsilon_{t-2} + \varepsilon_t
\end{equation}
where $B$ is the backshift operator, $\varepsilon_t \sim \mathcal{N}(0, \sigma^2)$
is white noise, $\theta_1$ and $\theta_2$ are the MA coefficients estimated
from the data, and $P_t$ is the weekly mean price.  The selected AIC is
170.2 and BIC is 171.1.

\textbf{Figure 15: SARIMAX Forecast} overlays the ARIMA(0,1,2) in-sample
fitted values and June 2019 out-of-sample forecast against the actual weekly
mean prices.  Visually, the model captures the broad upward trend from March
to May but over-smooths the June peak.  The forecast $R^2$ is $-0.3701$,
confirming that the 18-observation model cannot adequately predict the
June seasonal spike.

\begin{figure}[htbp]
\centering
\includegraphics[width=0.85\textwidth]{15_sarimax_forecast.png}
\caption{Figure 15 --- SARIMAX ARIMA(0,1,2) weekly forecast vs.\ actual prices. The model captures the March--May trend but over-smooths the June uptick; out-of-sample $R^2 = -0.3701$.}
\label{fig:f15}
\end{figure}

\subsection{Stage 7: Residual Diagnostics --- Ljung-Box Test (Figure 17)}

\textbf{Purpose}: Validate whether the SARIMAX model's residuals are white noise.

The Ljung-Box test evaluates the null hypothesis that the first $m$ sample
autocorrelations of the residuals are jointly zero:
\begin{equation}
Q(m) = n(n+2)\sum_{k=1}^{m}\frac{\hat{\rho}_k^2}{n-k} \sim \chi^2(m)
\end{equation}
where $\hat{\rho}_k$ are the sample autocorrelations of model residuals and
$n$ is the number of observations.

\textbf{Figure 17: Residual Diagnostics} displays: (a) the ACF of model
residuals, (b) a Q-Q plot of residuals versus a normal distribution, and
(c) the time plot of residuals.  The ACF of residuals shows spikes exceeding
the 95\% confidence band at lags 10 and 20, visible in the figure.
Numerically:

\begin{itemize}
    \item Lag 5: $Q = 7.068$, $p = 0.2156$ --- \textbf{fail to reject $H_0$}
    (white noise; good local fit)
    \item Lag 10: $Q = 21.029$, $p = 0.0209$ --- \textbf{reject $H_0$}
    (significant residual autocorrelation at medium lag)
    \item Lag 20: $Q = 33.075$, $p = 0.0331$ --- \textbf{reject $H_0$}
    (significant residual autocorrelation at long lag)
\end{itemize}

The short-lag white noise confirms that the MA(2) structure captures immediate
price persistence well.  The medium- and long-lag significance indicates
that there are periodic price structures at 10--20 week horizons that the
ARIMA(0,1,2) cannot capture---consistent with annual seasonality patterns
that require more data than 18 observations to model.

\begin{figure}[htbp]
\centering
\includegraphics[width=\textwidth]{17_residual_diagnostics.png}
\caption{Figure 17 --- SARIMAX residual diagnostics. ACF of residuals shows spikes at lags 10 and 20 (significant autocorrelation); Q-Q plot indicates approximately normal residuals; Ljung-Box $p < 0.05$ at lags 10 and 20.}
\label{fig:f17}
\end{figure}

\subsection{Stage 8: Hybrid LightGBM + SARIMAX (Figures 18--19)}
\label{sec:hybrid_method}

\textbf{Purpose}: Combine the trend signal from SARIMAX with the feature-rich
predictions of LightGBM to outperform either approach alone under the temporal
split.

\textbf{Architecture}:
\begin{enumerate}
    \item Fit SARIMAX ARIMA(0,1,2) on the 18 weekly mean price observations.
    \item Extract the \textbf{in-sample fitted weekly price}
    $\hat{P}_{\text{week}(t)}$ for every week in the training period.
    \item For each individual flight record, look up the fitted price for
    its journey week and append it as a new feature column:
    \texttt{SARIMAX\_weekly\_trend}.
    \item Retrain LightGBM on the temporal training set (March--May) with
    this 46th feature added.
    \item Evaluate on the June test set (using SARIMAX June forecast values
    as the trend feature for test records).
\end{enumerate}

\textbf{Intuition}: LightGBM can model which airline, route, and stop count
cause a ticket to cost more or less than the population average, but it
cannot predict the absolute price level in June if June prices are
systematically higher than May.  The SARIMAX fitted value provides the
macro-level ``where are prices this week?'' signal, while LightGBM provides
the micro-level ``how does this specific ticket deviate from the weekly average?''
signal.

\textbf{Figure 18: Hybrid Scatter} overlays actual vs.\ predicted prices
for the hybrid model on the June 2019 test set, compared to the standalone
temporal LightGBM.  The hybrid scatter cloud is noticeably tighter around
the diagonal (perfect prediction line), visually confirming the $R^2$
improvement.

\begin{figure}[htbp]
\centering
\includegraphics[width=0.75\textwidth]{18_hybrid_scatter.png}
\caption{Figure 18 --- Actual vs.\ predicted scatter for the hybrid LightGBM+SARIMAX model (temporal split, June 2019 test set). Points cluster more tightly around the diagonal compared to the standalone temporal LightGBM ($R^2 = 0.7629$ vs.\ 0.5351).}
\label{fig:f18}
\end{figure}

\textbf{Figure 19: Full Model Comparison} is a comprehensive bar chart
comparing all five baselines, temporal LightGBM, SARIMAX, and the hybrid
model across $R^2$, MAE, and RMSE.  The hybrid's bar in the $R^2$ chart
sits between the temporal baseline (0.5351) and the random-split baselines
($\geq 0.79$), illustrating the substantial recovery achieved by the
SARIMAX trend feature.

\begin{figure}[htbp]
\centering
\includegraphics[width=\textwidth]{19_full_model_comparison.png}
\caption{Figure 19 --- Full model comparison bar charts ($R^2$, MAE, RMSE) across all 11 evaluated configurations. The hybrid model sits clearly above the temporal LightGBM baseline and approaches the random-split results.}
\label{fig:f19}
\end{figure}

\subsection{Stage 9: Optuna Hyperparameter Tuning (Figure 11)}

\textbf{Purpose}: Find better hyperparameter configurations for the three
strongest baseline models to improve generalisation.

\textbf{Optimiser}: Optuna with Tree-structured Parzen Estimator (TPE)
sampling~\cite{akiba2019}.  TPE builds probabilistic models of the
objective function from previous trial results to propose the most promising
next trial configuration.  This is substantially more efficient than random
search and avoids the curse of dimensionality that makes grid search
impractical for high-dimensional search spaces.

\textbf{Objective}: 5-fold cross-validated $R^2$ on the training set.
Evaluating on held-out data within the cross-validation prevents overfitting
the hyperparameters to a single train/test split.

\textbf{Trials}: 50 per model.  At 5-fold CV per trial, this amounts to
250 full model fits per algorithm---computationally intensive but feasible
with LightGBM and XGBoost's efficient implementations.

\textbf{Search spaces and selected parameters} for each model:

\begin{itemize}
    \item \textbf{Random Forest}: searched \texttt{n\_estimators} [100, 1000],
    \texttt{max\_depth} [5, 30 or None], \texttt{min\_samples\_split} [2, 20],
    \texttt{min\_samples\_leaf} [1, 10].

    \item \textbf{XGBoost}: searched \texttt{n\_estimators} [100, 1000],
    \texttt{max\_depth} [3, 12], \texttt{learning\_rate} [$10^{-3}$, 0.3]
    (log scale), \texttt{subsample} [0.5, 1.0], \texttt{colsample\_bytree}
    [0.5, 1.0], \texttt{reg\_alpha} and \texttt{reg\_lambda} [$10^{-8}$, 10.0]
    (log scale).  The selected configuration used a lower learning rate
    ($\approx 0.05$) and deeper trees (\texttt{max\_depth} $\approx 8$) than
    defaults.

    \item \textbf{LightGBM}: searched \texttt{n\_estimators} [100, 1000],
    \texttt{max\_depth} [3, 12], \texttt{learning\_rate} [$10^{-3}$, 0.3]
    (log scale), \texttt{num\_leaves} [20, 300], \texttt{min\_child\_samples}
    [5, 100], \texttt{subsample} [0.5, 1.0], \texttt{colsample\_bytree}
    [0.5, 1.0].
\end{itemize}

\textbf{Figure 11: Tuned Model Comparison} plots the same three-metric bar
chart as Figure 10 but for the tuned models.  XGBoost rises to the top
position after tuning, overtaking LightGBM in $R^2$ (0.8400 vs.\ 0.8371)
and achieving the lowest MAE (Rs.\,1{,}162).  The visual gap between the
tuned models is narrow, confirming that all three are competitive after
optimisation.

\begin{figure}[htbp]
\centering
\includegraphics[width=\textwidth]{11_tuned_model_comparison.png}
\caption{Figure 11 --- Tuned model comparison (50 Optuna trials each). XGBoost overtakes LightGBM after tuning; all three tuned models achieve $R^2 > 0.81$.}
\label{fig:f11}
\end{figure}

\subsection{Stage 10: Test Set Prediction (Figure 20)}

\textbf{Purpose}: Apply the best tuned model to the held-out test set and
visualise prediction quality across the full price range.

The tuned XGBoost model is applied to the 2{,}093 test records.
\textbf{Figure 20: Test Predictions} is a time-ordered scatter plot of
actual vs.\ predicted prices.  The plot shows that:
\begin{itemize}
    \item Predictions in the Rs.\,3{,}000--Rs.\,15{,}000 range are tight,
    confirming accurate modelling of the majority of the price distribution.
    \item At the extremes (fares above Rs.\,30{,}000), predictions
    systematically underestimate---a consequence of the right tail having
    fewer training examples and RMSE-minimising loss applying weaker
    penalty to underestimation of extreme values.
    \item The model shows no systematic directional bias in the central
    price range, indicating well-calibrated predictions for the typical
    passenger.
\end{itemize}

\begin{figure}[htbp]
\centering
\includegraphics[width=0.8\textwidth]{20_test_predictions.png}
\caption{Figure 20 --- Tuned XGBoost predictions on the held-out test set. Accurate across the Rs.\,3{,}000--15{,}000 range; systematic under-prediction for fares above Rs.\,30{,}000 due to sparse training examples in the right tail.}
\label{fig:f20}
\end{figure}

\subsection{Stage 11: Route-Level Fare Forecasting (Figures 21--23)}

\textbf{Purpose}: Generate actionable airline-stratified fare estimates
for two high-traffic domestic corridors.

\textbf{Construction of prediction feature vectors}: For each airline on
each corridor, a synthetic booking record is created using:
\begin{itemize}
    \item Observed average values for that airline/route combination
    (typical number of stops, average departure hour, typical flight duration).
    \item Journey date swept across all June 2019 days (1--30).
    \item SARIMAX weekly trend feature set to the June forecast value.
    \item All other features set to the airline-route-typical values
    observed in training data.
\end{itemize}

The tuned XGBoost model (best overall performance) is used for inference.

\textbf{Figure 21: Jet Airways Delhi--Cochin Forecast} shows predicted daily
fares for Jet Airways on the Delhi--Cochin corridor throughout June 2019.
The line is relatively flat (Rs.\,10{,}577--Rs.\,11{,}600) with a slight
mid-month dip and end-of-month increase.  The stability reflects Jet Airways'
full-service pricing policy: unlike LCCs that vary prices more aggressively,
full-service carriers maintain narrower pricing bands on specific itineraries.

\begin{figure}[htbp]
\centering
\includegraphics[width=0.8\textwidth]{21_forecast_jetairways_del_cok.png}
\caption{Figure 21 --- Jet Airways Delhi--Cochin predicted daily fares for June 2019 (Rs.\,10{,}577--Rs.\,11{,}600). Stable premium pricing with narrow day-to-day variation.}
\label{fig:f21}
\end{figure}

\textbf{Figure 22: IndiGo Bangalore--Delhi Forecast} shows a nearly flat
daily fare profile (Rs.\,4{,}567--Rs.\,4{,}714) for IndiGo on the
Bangalore--Delhi corridor.  The coefficient of variation is under 3\%,
consistent with IndiGo's historical approach of stable, high-volume low-cost
pricing on its densest routes.

\begin{figure}[htbp]
\centering
\includegraphics[width=0.8\textwidth]{22_forecast_indigo_blr_del.png}
\caption{Figure 22 --- IndiGo Bangalore--Delhi predicted daily fares for June 2019 (mean Rs.\,4{,}614, CV $< 3\%$). Near-flat profile consistent with high-frequency LCC pricing stability.}
\label{fig:f22}
\end{figure}

\textbf{Figure 23: All-Airline Comparison (Delhi--Cochin)} overlays predicted
June 2019 fares for all eight airlines serving the Delhi--Cochin corridor
on a single chart.  The two-tier pricing structure is clearly visible:
\begin{itemize}
    \item \textbf{Premium tier}: Multiple Carriers (Rs.\,12{,}483), Air India
    (Rs.\,12{,}408), Jet Airways (Rs.\,11{,}255)
    \item \textbf{Budget tier}: SpiceJet (Rs.\,7{,}887), Air Asia (Rs.\,7{,}027),
    IndiGo (Rs.\,6{,}833), Vistara (Rs.\,6{,}748), GoAir (Rs.\,6{,}255)
\end{itemize}
The gap between tiers is approximately Rs.\,4{,}000--5{,}000 (40--50\% of the
budget tier fare), quantifying the consumer cost of choosing a full-service
airline on this corridor.

\begin{figure}[htbp]
\centering
\includegraphics[width=\textwidth]{23_forecast_comparison_del_cok.png}
\caption{Figure 23 --- All-airline June 2019 fare forecast on the Delhi--Cochin corridor. Clear two-tier structure: full-service carriers (Rs.\,11{,}000--12{,}500) vs.\ budget LCCs (Rs.\,6{,}255--7{,}887).}
\label{fig:f23}
\end{figure}

\subsection{Stage 12: Actual vs.\ Predicted Visualisation (Figures 24--25)}

These figures provide diagnostic scatter plots for the two most important
models.

% ===========================================================================
\section{Experimental Results}
\label{sec:results}
% ===========================================================================

\subsection{Baseline Model Comparison}

Table~\ref{tab:baseline} reports results for all five baseline models on
the random 20\% hold-out set ($n = 2{,}093$).

\begin{table}[h]
\caption{Baseline Model Comparison --- Random 80/20 Split}
\label{tab:baseline}
\centering
\begin{tabular}{lrrrr}
\toprule
\textbf{Model} & \textbf{Train $R^2$} & \textbf{Test $R^2$} & \textbf{MAE (Rs.)} & \textbf{RMSE (Rs.)} \\
\midrule
Linear Regression      & 0.5925 & 0.6037 & 1{,}997 & 2{,}875 \\
Decision Tree          & 0.9697 & 0.6277 & 1{,}418 & 2{,}786 \\
Random Forest          & 0.9510 & 0.8052 & 1{,}187 & 2{,}015 \\
XGBoost                & 0.9439 & 0.7908 & 1{,}205 & 2{,}088 \\
\textbf{LightGBM}      & \textbf{0.8796} & \textbf{0.8275} & 1{,}261 & \textbf{1{,}896} \\
\bottomrule
\end{tabular}
\end{table}

\textbf{Key observations}:

\begin{itemize}
    \item \textbf{Linear Regression} ($R^2 = 0.6037$): Approximately 40\% of
    price variance remains unexplained by any linear combination of features.
    The relatively high MAE of Rs.\,1{,}997 (USD 28.36) makes this model
    insufficient for practical budgeting applications.

    \item \textbf{Decision Tree} ($R^2 = 0.6277$): Only marginally better than
    linear regression on the test set despite near-perfect training fit
    (0.9697).  The train-test gap of 0.342 points is the largest among all
    models and definitively confirms that an unpruned decision tree memorises
    training data rather than generalising.  RMSE of Rs.\,2{,}786 is slightly
    better than linear regression because the tree captures non-linear splits
    but collapses badly on unseen combinations.

    \item \textbf{Random Forest} ($R^2 = 0.8052$): Bagging over 100 trees
    effectively reduces variance, producing the best Random Forest result
    without any tuning.  Train-test gap of 0.145 is reasonable.  MAE of
    Rs.\,1{,}187 makes this the lowest-MAE baseline model.

    \item \textbf{XGBoost} ($R^2 = 0.7908$, MAE = Rs.\,1{,}205$): Slightly
    below Random Forest under default parameters.  XGBoost's regularisation
    prevents the extreme overfitting seen in the Decision Tree (train $R^2$
    = 0.9439 vs.\ 0.8052 for Random Forest), but its default learning rate
    and tree depth are not optimal for this dataset without tuning.

    \item \textbf{LightGBM} ($R^2 = 0.8275$, RMSE = Rs.\,1{,}896$): Best
    cross-sectional performance.  LightGBM's leaf-wise (best-first) tree
    growth strategy allocates more splits to high-gain nodes, which is
    particularly advantageous for right-skewed distributions like fare prices.
    The train-test gap of 0.052 is the smallest among ensemble models,
    indicating the best generalisation of the default configurations.
\end{itemize}

\subsection{Temporal Split Analysis and Figure 24}

\textbf{Figure 24: Temporal Actual vs.\ Predicted (Daily Mean)} shows the
daily mean actual fare versus the LightGBM predicted daily mean fare for
all June 2019 test records, grouped by journey day.  This plot makes the
temporal distribution shift visible:

\begin{itemize}
    \item Days 1--7 of June: predictions track actuals closely; the model
    generalises well from late-May patterns.
    \item Days 8--16: predictions underestimate actual prices by Rs.\,500--1{,}500
    on several days, corresponding to the pre-monsoon travel surge (weekend
    of June 14--16 is a public holiday weekend in several Indian states).
    \item Days 17--27: predictions partially recover as prices normalise.
    \item The shaded fill between actual and predicted lines (light grey)
    makes the systematic underestimation on high-demand days visually clear.
\end{itemize}

This plot directly explains why the temporal split yields a lower $R^2$
(0.5351) than the random split (0.8275): the model cannot anticipate
event-driven demand spikes that did not occur in the training period.

\begin{figure}[htbp]
\centering
\includegraphics[width=\textwidth]{24_temporal_actual_vs_predicted.png}
\caption{Figure 24 --- Daily mean actual vs.\ predicted fare for LightGBM under the temporal split (June 2019 test set). Systematic underestimation on days 8--16 reveals the June demand spike the model cannot anticipate from March--May training data.}
\label{fig:f24}
\end{figure}

\subsection{SARIMAX Results}

On 18 weekly observations, auto-ARIMA selected ARIMA(0,1,2) with:
\begin{itemize}
    \item AIC = 170.2 (model fit quality; lower is better)
    \item BIC = 171.1 (penalises parameters more heavily; similar to AIC
    confirms the 2-parameter MA model is appropriate)
    \item Forecast $R^2 = -0.3701$ (June weekly mean predictions are worse
    than predicting the grand mean)
    \item Forecast MAE $= \text{Rs.}\ 1{,}282$ (weekly aggregate comparison;
    not directly comparable to flight-level ML MAE)
\end{itemize}

The negative $R^2$ does not mean the SARIMAX model is useless---it means
that for the specific task of predicting week-by-week absolute price levels
in June from only 12 historical weekly observations (March through May), the
model over-smooths and misses the June uptick.  However, the in-sample fitted
values (which use all 18 observations) are well-calibrated and provide the
trend signal exploited by the hybrid model.

\subsection{Ljung-Box Residual Diagnostics}

\begin{table}[h]
\caption{Ljung-Box Test Results --- SARIMAX ARIMA(0,1,2) Residuals}
\label{tab:ljungbox}
\centering
\begin{tabular}{crrl}
\toprule
\textbf{Lag} & \textbf{Q-statistic} & \textbf{p-value} & \textbf{Interpretation} \\
\midrule
5  & 7.068  & 0.2156 & White noise at short lag (PASS) \\
10 & 21.029 & 0.0209 & Residual autocorrelation at lag 10 (FAIL) \\
20 & 33.075 & 0.0331 & Residual autocorrelation at lag 20 (FAIL) \\
\bottomrule
\end{tabular}
\end{table}

The lag-5 pass confirms that the MA(2) structure adequately captures
immediate price persistence.  The lag-10 and lag-20 failures indicate
periodic patterns at 10-week and 20-week intervals that the model cannot
capture with only 18 observations---essentially, these are signatures of
annual seasonality (Indian travel peaks in festival season, 30--40 weeks
away from the data window) that the model cannot identify.

\subsection{Hybrid Model Results}

\begin{table}[h]
\caption{Hybrid LightGBM + SARIMAX vs.\ Temporal LightGBM}
\label{tab:hybrid}
\centering
\begin{tabular}{lrr}
\toprule
\textbf{Model} & \textbf{Test $R^2$} & \textbf{MAE (Rs.)} \\
\midrule
LightGBM (temporal, no TS feature) & 0.5351 & 1{,}758 \\
LightGBM + SARIMAX (hybrid)         & 0.7629 & 1{,}392 \\
\midrule
\textbf{Absolute improvement}       & \textbf{+0.2278} & \textbf{$-$366} \\
\textbf{Relative improvement (MAE)} & --- & \textbf{$-$20.8\%} \\
\bottomrule
\end{tabular}
\end{table}

Adding the single SARIMAX weekly trend feature improves $R^2$ by 22.78
percentage points (from 0.5351 to 0.7629) and reduces MAE by Rs.\,366
(20.8\%).  This is a substantial gain from a single feature, confirming
the hybrid architecture hypothesis.  The hybrid $R^2 = 0.7629$ approaches
the random-split LightGBM baseline (0.8275)---most of the temporal
performance penalty has been recovered.

\textbf{Why is the gain so large?}  The SARIMAX trend feature directly
encodes the macro price level shift from May to June.  Without it,
LightGBM's training data (March--May) has a mean price of approximately
Rs.\,8{,}985, while the June test set mean is Rs.\,9{,}610.  The model
therefore anchors predictions to the lower training mean.  The SARIMAX
feature tells the model: ``this flight is in a week where prices are
roughly Rs.\,9{,}400'', correcting the systematic downward bias.

\subsection{Hyperparameter Tuning Results}

\begin{table}[h]
\caption{Optuna-Tuned Model Performance (50 Trials, Random 80/20 Split)}
\label{tab:tuned}
\centering
\begin{tabular}{lrrrrr}
\toprule
\textbf{Model} & \textbf{CV $R^2$} & \textbf{Test $R^2$} & \textbf{MAE (Rs.)} & \textbf{RMSE (Rs.)} & \textbf{MAE (USD)} \\
\midrule
Random Forest (tuned) & 0.7959 & 0.8188 & 1{,}233 & 1{,}944 & \$17.52 \\
XGBoost (tuned)       & 0.8188 & \textbf{0.8400} & \textbf{1{,}162} & \textbf{1{,}827} & \textbf{\$16.50} \\
LightGBM (tuned)      & 0.8179 & 0.8371 & 1{,}182 & 1{,}843 & \$16.78 \\
\bottomrule
\end{tabular}
\end{table}

Tuning gains relative to default baselines:
\begin{itemize}
    \item Random Forest: $+0.013$ $R^2$ points (from 0.8052 to 0.8188)
    \item XGBoost: $+0.049$ $R^2$ points (from 0.7908 to 0.8400) --- largest gain
    \item LightGBM: $+0.010$ $R^2$ points (from 0.8275 to 0.8371) --- smallest gain
\end{itemize}

XGBoost benefits most from tuning because its default learning rate (0.3)
is too aggressive for this dataset, causing premature convergence to a
suboptimal solution.  The Optuna search found that a lower learning rate
($\approx 0.05$) paired with more trees and regularisation achieves better
generalisation.

LightGBM benefits least because its default configuration is already
well-calibrated for structured tabular regression; the leaf-wise growth
strategy naturally adapts to the data's right-skewed price distribution
without the aggressive learning rate penalty that affects XGBoost.

Tuned Random Forest ($R^2 = 0.8188$) underperforms default LightGBM
(0.8275), confirming that Random Forest's architecture---which uses symmetric
(depth-wise) tree growth and a fixed maximum depth---has less headroom for
improvement on this particular feature distribution than the gradient boosting
approaches.

\subsection{Figure 25: Actual vs.\ Predicted --- Tuned XGBoost}

\textbf{Figure 25} presents a two-panel diagnostic for the best model
(Tuned XGBoost on the random hold-out set):

\textbf{Left panel --- Scatter plot (Actual vs.\ Predicted)}:
The plot title shows $R^2 = 0.8400$ and MAE $= $ Rs.\,1{,}162.  The
majority of points cluster tightly around the diagonal (perfect prediction
line, shown in red dashes).  Notable observations:
\begin{itemize}
    \item Points below the diagonal (under-predictions): concentrated at
    actual prices above Rs.\,20{,}000, corresponding to premium Jet Airways
    and Air India tickets.  These are systematically under-predicted because
    the model has seen fewer extreme-price examples.
    \item Points above the diagonal (over-predictions): sparse, occurring
    mainly in the Rs.\,5{,}000--10{,}000 range, where the model occasionally
    predicts a higher tier than the actual ticket.
    \item The central mass (Rs.\,3{,}000--15{,}000) is well-concentrated
    around the diagonal, confirming reliable predictions for typical consumers.
\end{itemize}

\textbf{Right panel --- Residual distribution histogram}:
The residuals (Predicted $-$ Actual) are plotted as a histogram with a red
dashed line at zero.  The distribution is approximately normal and centred
near zero, indicating no systematic bias.  The distribution has light tails
extending to $\pm$Rs.\,10{,}000, consistent with the RMSE/MAE ratio analysis
in Section~\ref{sec:metrics}: most errors are small (MAE = Rs.\,1{,}162)
but a minority of predictions are off by Rs.\,5{,}000--10{,}000.

\begin{figure}[htbp]
\centering
\includegraphics[width=\textwidth]{25_actual_vs_predicted_tuned.png}
\caption{Figure 25 --- Tuned XGBoost diagnostic ($R^2 = 0.8400$, MAE = Rs.\,1{,}162). \textit{Left}: scatter of actual vs.\ predicted prices; points cluster tightly around the diagonal in the Rs.\,3{,}000--15{,}000 range with some underestimation at the premium tail. \textit{Right}: residual histogram -- approximately normal, centred at zero, confirming no systematic bias.}
\label{fig:f25}
\end{figure}

\subsection{Complete Model Comparison Summary}

Table~\ref{tab:full_comparison} consolidates all 11 evaluated models.

\begin{table}[h]
\caption{Complete Model Evaluation Summary --- All Models}
\label{tab:full_comparison}
\centering
\begin{tabular}{llrrr}
\toprule
\textbf{Category} & \textbf{Model} & \textbf{$R^2$} & \textbf{MAE (Rs.)} & \textbf{RMSE (Rs.)} \\
\midrule
\multirow{5}{*}{Baseline (random)}
  & Linear Regression      & 0.6037 & 1{,}997 & 2{,}875 \\
  & Decision Tree          & 0.6277 & 1{,}418 & 2{,}786 \\
  & Random Forest          & 0.8052 & 1{,}187 & 2{,}015 \\
  & XGBoost                & 0.7908 & 1{,}205 & 2{,}088 \\
  & LightGBM               & 0.8275 & 1{,}261 & 1{,}896 \\
\midrule
\multirow{3}{*}{Temporal split}
  & LightGBM (temporal)       & 0.5351 & 1{,}758 & --- \\
  & SARIMAX ARIMA(0,1,2)       & $-$0.3701 & 1{,}282 & --- \\
  & LightGBM + SARIMAX (hybrid) & 0.7629 & 1{,}392 & --- \\
\midrule
\multirow{3}{*}{Tuned (random)}
  & Random Forest (tuned)  & 0.8188 & 1{,}233 & 1{,}944 \\
  & \textbf{XGBoost (tuned)} & \textbf{0.8400} & \textbf{1{,}162} & \textbf{1{,}827} \\
  & LightGBM (tuned)       & 0.8371 & 1{,}182 & 1{,}843 \\
\bottomrule
\end{tabular}
\end{table}

\subsection{Route-Level Forecasting Results}

\begin{table}[h]
\caption{Predicted June 2019 Fares: Delhi $\rightarrow$ Cochin, All Airlines}
\label{tab:del_cok}
\centering
\begin{tabular}{lrr}
\toprule
\textbf{Airline} & \textbf{Mean Predicted Fare (Rs.)} & \textbf{Tier} \\
\midrule
Multiple Carriers & 12{,}483 & Full / Mixed \\
Air India         & 12{,}408 & Full \\
Jet Airways       & 11{,}255 & Full \\
SpiceJet          &  7{,}887 & Budget \\
Air Asia          &  7{,}027 & Budget \\
IndiGo            &  6{,}833 & Budget \\
Vistara           &  6{,}748 & Budget \\
GoAir             &  6{,}255 & Budget \\
\midrule
\textbf{Full-service average} & \textbf{12{,}049} & \\
\textbf{Budget average}       & \textbf{6{,}950}  & \\
\textbf{Full / Budget ratio}  & \textbf{1.73:1}   & \\
\bottomrule
\end{tabular}
\end{table}

Jet Airways' predicted fares range from Rs.\,10{,}577 to Rs.\,11{,}600
across the June calendar, while IndiGo on the same corridor averages
Rs.\,6{,}833.  The Rs.\,4{,}422 average gap (39\% of IndiGo's fare)
represents the premium consumers pay for full-service features (meals,
checked baggage, wider seats, airport lounge access) on this route.

% ===========================================================================
\section{Discussion}
\label{sec:discussion}
% ===========================================================================

\subsection{Key Findings}

\textbf{Finding 1 --- Non-linearity is fundamental to fare structure.}
Linear Regression explains only 60.4\% of fare variance, while LightGBM
explains 82.8\%.  The 22.4 percentage-point gap is large and confirms that
fare determination involves complex multiplicative interactions (e.g., a
Jet Airways ticket on a non-stop long-haul route costs disproportionately
more than the sum of the ``Jet Airways premium'' and the ``non-stop premium'').
Any practical fare estimation tool must use a non-linear model family.

\textbf{Finding 2 --- Temporal evaluation substantially changes conclusions.}
Under random splitting, LightGBM achieves $R^2 = 0.8275$.  Under temporal
splitting (the realistic deployment scenario), performance drops to 0.5351.
This 29-point gap represents the cost of ``looking into the future'' during
training: a randomly split model has seen June-priced tickets during
training, making June prediction trivially easier.  Studies that report
only random-split metrics significantly overstate deployment performance.

\textbf{Finding 3 --- The hybrid architecture substantially mitigates the
temporal performance penalty.}  The SARIMAX weekly trend feature recovers
22.78 percentage points of the 29-point temporal penalty.  The residual
7-point gap ($0.8275 - 0.7629 = 0.065$) represents the irreducible cost
of genuine distribution shift between spring and early-summer booking
patterns---demand spikes from public holidays and school holidays that the
model cannot anticipate from historical data alone.

\textbf{Finding 4 --- Bayesian hyperparameter optimisation yields consistent
but modest gains.}  Across the three tuned models, the average $R^2$
improvement is $+0.024$ over defaults.  The largest gain is for XGBoost
($+0.049$).  These gains are real and worth obtaining in production, but
they do not change the ranking of modelling approaches: LightGBM and
XGBoost were already the best models before tuning.

\textbf{Finding 5 --- Number of stops, not airline brand, is the primary
pricing signal.}  With importance of 23.38\%, Total\_Stops outweighs
Airline\_Tier (14.53\%) and the Jet Airways dummy (9.29\%).  This reveals
that \textit{how} a passenger travels (direct vs.\ connecting) matters more
than \textit{who} they fly with in terms of fare determination.  Airlines
use stop-count as a proxy for route difficulty and schedule optimisation:
non-stop flights on long-haul routes require premium aircraft, dedicated
crew scheduling, and command higher demand, all of which inflate fares.

\subsection{Limitations}

\textbf{Short temporal window}: Four months (March--June 2019) do not capture
annual seasonality (Diwali in October, Christmas in December, summer school
holidays in May).  A model trained exclusively on spring data may not transfer
to autumn or winter without retraining.

\textbf{Static booking snapshot}: Each record represents a single-point-in-time
fare, not the full trajectory of price changes from 30 days before departure
to the day of departure.  A dynamic fare prediction model would require
time-stamped fare sequences per itinerary---a richer data collection challenge.

\textbf{Jet Airways collapse anomaly}: Jet Airways suspended operations on
April 17, 2019, mid-dataset.  Records from April and early May likely reflect
distressed pricing anomalies (last-minute premium fares as passengers scrambled,
competitor capacity absorption, price spikes on previously Jet-served routes).
These anomalies could inflate prediction error for Jet Airways tickets and
contaminate the temporal signal around April.

\textbf{No external regressors}: Jet fuel prices, airport congestion indices,
public holiday calendars, and competitor seat availability are absent from the
feature set.  SARIMAX with well-chosen exogenous variables and more data would
likely produce a stronger time-series component.

\textbf{Aggregation loss in time-series models}: Any classical TS model
requires aggregating individual flights to a single daily or weekly price,
discarding the airline, route, stop, and departure-time features that account
for over 80\% of predictive importance.  This fundamental mismatch is why
TS models alone cannot compete with flight-level ML approaches.

\subsection{Practical Implications}

\textbf{For consumers}: The tuned XGBoost model provides fare estimates
with MAE $=$ Rs.\,1{,}162 (USD 16.50).  For a typical domestic booking
in the Rs.\,5{,}000--15{,}000 range, this represents an error of 8--23\%,
which is useful for budget planning (``the flight will cost approximately
Rs.\,8{,}000 $\pm$ Rs.\,1{,}200'') but not precise enough for automated
purchase-timing decisions.

\textbf{For airlines}: The dominance of stop-count (23.38\%) and airline tier
(14.53\%) over temporal booking signals (Journey\_Day: 10.81\%, Journey\_Month:
1.93\%) implies that consumers are primarily price-differentiated by routing
complexity and service tier, not by when they book.  Revenue management
systems should therefore focus on route design and tier differentiation rather
than pure yield-management timing.

\textbf{For regulators}: The 1.73:1 full-service to budget fare ratio on
the Delhi--Cochin corridor quantifies the competitive pricing structure.
If Jet Airways' shutdown (April 2019) removed a full-service option, the
model predicts that remaining full-service alternatives (Air India) would
absorb demand at even higher prices, while budget carriers would experience
increased demand pressure.

% ===========================================================================
\section{Conclusion}
\label{sec:conclusion}
% ===========================================================================

This paper presented a comprehensive, end-to-end machine learning and time
series pipeline for Indian domestic flight fare prediction.  Starting from
10{,}683 raw booking records, we delivered 11 trained models, 25 analytical
visualisations, and 10 structured output reports, covering the full spectrum
from basic EDA to Bayesian-optimised ensemble models and a hybrid
ML+time-series architecture.

The tuned XGBoost regressor achieves $R^2 = 0.8400$ and MAE $=$ Rs.\,1{,}162
on a held-out random test set, representing the best cross-sectional
performance.  Under a rigorous temporal split (March--May training, June
testing), the hybrid LightGBM+SARIMAX model achieves $R^2 = 0.7629$,
recovering 22.78 percentage points of the 29-point temporal performance penalty
incurred by the standalone ML model.

We identified and corrected a systematic duration-parsing error ($10\times$
magnitude corruption) present in widely used public implementations, formally
evaluated four time-series paradigms for applicability to fare data (GARCH and
SVAR ruled out; SARIMA and SARIMAX limited to supplementary trend roles), and
generated airline-stratified forecasts revealing a 1.73:1 full-service to
budget pricing divide on India's busiest domestic corridor.

\textbf{Future work} will pursue:
\begin{enumerate}
    \item Extension to $\geq 2$ years of booking data to enable reliable annual
    seasonality modelling with SARIMA ($s = 365$).
    \item Dynamic fare trajectory modelling: tracking how a specific ticket's
    price changes from 60 days before departure to the booking date.
    \item Incorporation of exogenous regressors into SARIMAX: jet fuel prices,
    national holiday calendars, airline seat load factors.
    \item Deep learning alternatives: TabNet and FT-Transformer have shown
    competitive performance on structured tabular regression with sufficient data.
    \item Causal modelling of competitor pricing: using SVAR once multi-airline
    route-level panel data becomes available to model demand spillovers from
    Jet Airways' collapse.
    \item Real-time deployment: a web API that accepts a proposed itinerary and
    returns a fare estimate with confidence interval, based on retrained weekly
    models.
\end{enumerate}

% ===========================================================================
\begin{thebibliography}{00}

\bibitem{dgca2019}
Directorate General of Civil Aviation (DGCA), India,
``Annual Report 2018--19,''
Ministry of Civil Aviation, Government of India, New Delhi, 2019.

\bibitem{belobaba2009}
P.~Belobaba, A.~Odoni, and C.~Barnhart, Eds.,
\textit{The Global Airline Industry}.
Chichester, UK: John Wiley \& Sons, 2009.

\bibitem{grover2015}
P.~Grover,
``Predicting airline ticket prices,''
\textit{CS229 Machine Learning Project Reports},
Stanford University, Stanford, CA, USA, 2015.

\bibitem{janssen2019}
T.~Janssen, A.~Dirkzwager, and O.~Postma,
``Airline price prediction with machine learning,''
in \textit{Proc. 1st Int. Conf. Data Science, E-Learning and Information
Systems (ICDSEIS)}, Madrid, Spain, Oct. 2019, pp. 1--5.

\bibitem{raghav2020}
K.~Raghav and A.~Sinha,
``Flight fare prediction using machine learning,''
\textit{Int. J. Recent Technol. Eng.},
vol.~8, no.~3, pp. 3564--3568, 2019.

\bibitem{tziridis2017}
K.~Tziridis, T.~Kalampokas, G.~A.~Papakostas, and K.~I.~Diamantaras,
``Airfare prices prediction using machine learning techniques,''
in \textit{Proc. 25th European Signal Processing Conf. (EUSIPCO)},
Kos Island, Greece, 2017, pp. 1036--1039.

\bibitem{etzioni2003}
O.~Etzioni, R.~Tuchinda, C.~A.~Knoblock, and A.~Yates,
``To buy or not to buy: Mining airfare data to minimize ticket purchase
price,''
in \textit{Proc. 9th ACM SIGKDD Int. Conf. Knowledge Discovery and Data
Mining (KDD)}, Washington, DC, USA, 2003, pp. 119--128.

\bibitem{grinsztajn2022}
L.~Grinsztajn, E.~Oyallon, and G.~Varoquaux,
``Why tree-based models still outperform deep learning on tabular data,''
in \textit{Advances in Neural Information Processing Systems (NeurIPS)},
vol.~35, 2022, pp. 507--520.

\bibitem{chen2016}
T.~Chen and C.~Guestrin,
``XGBoost: A scalable tree boosting system,''
in \textit{Proc. 22nd ACM SIGKDD Int. Conf. Knowledge Discovery and Data
Mining (KDD)}, San Francisco, CA, USA, 2016, pp. 785--794.

\bibitem{ke2017}
G.~Ke \textit{et al.},
``LightGBM: A highly efficient gradient boosting decision tree,''
in \textit{Advances in Neural Information Processing Systems (NIPS)},
vol.~30, 2017, pp. 3146--3154.

\bibitem{breiman2001}
L.~Breiman,
``Random forests,''
\textit{Machine Learning}, vol.~45, no.~1, pp. 5--32, 2001.

\bibitem{box2015}
G.~E.~P.~Box, G.~M.~Jenkins, G.~C.~Reinsel, and G.~M.~Ljung,
\textit{Time Series Analysis: Forecasting and Control}, 5th ed.
Hoboken, NJ, USA: Wiley, 2015.

\bibitem{shumway2017}
R.~H.~Shumway and D.~S.~Stoffer,
\textit{Time Series Analysis and Its Applications: With R Examples}, 4th ed.
Cham, Switzerland: Springer, 2017.

\bibitem{zhang2003}
G.~P.~Zhang,
``Time series forecasting using a hybrid ARIMA and neural network model,''
\textit{Neurocomputing}, vol.~50, pp. 159--175, Jan. 2003.

\bibitem{akiba2019}
T.~Akiba, S.~Sano, T.~Yanase, T.~Ohta, and M.~Koyama,
``Optuna: A next-generation hyperparameter optimization framework,''
in \textit{Proc. 25th ACM SIGKDD Int. Conf. Knowledge Discovery and Data
Mining (KDD)}, Anchorage, AK, USA, 2019, pp. 2623--2631.

\bibitem{bergstra2012}
J.~Bergstra and Y.~Bengio,
``Random search for hyper-parameter optimization,''
\textit{J. Mach. Learn. Res.}, vol.~13, no.~1, pp. 281--305, 2012.

\bibitem{bergstra2011}
J.~Bergstra, R.~Bardenet, Y.~Bengio, and B.~K\'{e}gl,
``Algorithms for hyper-parameter optimization,''
in \textit{Advances in Neural Information Processing Systems (NIPS)},
vol.~24, 2011, pp. 2546--2554.

\bibitem{grubb2001}
H.~Grubb and A.~Mason,
``Long lead-time forecasting of UK air passengers by Holt-Winters methods
with damped trend,''
\textit{Int. J. Forecasting}, vol.~17, no.~1, pp. 71--82, 2001.

\bibitem{grosche2007}
T.~Grosche, F.~Rothlauf, and A.~Heinzl,
``Gravity models for airline passenger volume estimation,''
\textit{J. Air Transport Management}, vol.~13, no.~4, pp. 175--183, 2007.

\end{thebibliography}

\end{document}
